\begin{abstract}
Language models and other transformers learn relational composition mechanisms to solve tasks such as tracking information about subjects (``green car'') to answer questions, or parallelising precomputing subpaths in graph-based pathfinding problems. There is a deep theoretical literature on vector composition methods, yet we lack empirical studies of what mechanisms transformers learn in practice. In particular, different composition methods affect sparse autoencoders (SAEs), a popular method for decomposing model activations, in different ways.
We present empirical evidence in a controlled attention-only transformer that ordered relational information can be encoded via a relative magnitude-based mechanism (by a weighted sum of vectors) rather than predicted direction-based mechanisms such as additive matrix binding. While absolute magnitude-based mechanisms have been reported for other architectures, to our knowledge this is the first controlled demonstration of a relative magnitude mechanism in attention-only transformers.
This result challenges the prevailing view in mechanistic interpretability research that transformer features can be viewed as binary and independent, and motivates a re-examination of these methods with respect to feature activation value and interactions between features at different values.
In future work, we will move beyond the toy setting and attempt to find evidence of these mechanisms in real-world transformers.
% avoid self-plagiarism with this disclosure:
Note that part of this work was previously published in a workshop setting.
\end{abstract}